% Source: source/普通高考/2015/2015重庆理.pdf, page 1.
% Type: program flowchart; pdfimages -list reports no embedded bitmap (vector line art).
% Grid evidence: original page rendered at 180 dpi; grid preview step=100, fine_step=50.
% Node-edge list: 开始→s=0,k=0→判断；判断(是)→k=k+2→s=s+1/k→loop merge→判断；
% 判断(否)→输出k→结束. Solid segments: all borders and directed connectors;
% dashed segments: none. Labels 是/否 are adjacent to their corresponding exits.
% Every node has local parindent=0pt and centered text; the loop arrow ends at
% the incoming main vertical edge and does not add a second arrow to that edge.
% Branch labels are attached directly to their outgoing edges with natural
% anchors; no hand-tuned coordinates or shifts are used for text placement.
\begin{tikzpicture}[
  x=1cm,
  y=0.9cm,
  >=Stealth,
  line width=0.45pt,
  every node/.style={font=\normalsize, inner sep=1pt, align=center, anchor=center,
    execute at begin node={\setlength{\parindent}{0pt}}},
  flow/.style={draw, line width=0.45pt},
  startstop/.style={flow, rounded corners=3pt, minimum width=1.2cm, minimum height=0.70cm},
  process/.style={flow, minimum height=0.72cm, align=center},
  io/.style={flow, trapezium, trapezium left angle=74, trapezium right angle=106,
    minimum height=0.72cm, align=center},
  decision/.style={flow, diamond, aspect=2.8, minimum width=3.35cm,
    minimum height=1.0cm, inner sep=1pt, align=center}
]
  \node[startstop] (start) at (0,10.2) {开始};
  \node[process, minimum width=2.30cm] (init) at (0,8.85) {$s=0, k=0$};
  \node[decision] (check) at (0,5.65) {};
  \node[io, minimum width=1.78cm] (output) at (0,3.95) {输出 $k$};
  \node[startstop] (finish) at (0,2.55) {结束};

  \node[process, minimum width=2.05cm] (increment) at (2.55,6.45) {$k=k+2$};
  \node[process, minimum width=2.55cm, minimum height=0.95cm] (sum) at (2.55,7.75)
    {$s=s+\dfrac{1}{k}$};
  \coordinate (loopJoin) at (0,8.35);
  \coordinate (branchRight) at (2.55,5.65);

  % Main path.
  \draw[->] (start.south) -- (init.north);
  \draw[->] (init.south) -- (check.north);
  \draw[->] (output.south) -- (finish.north);

  % Decision branches and loop.
  \draw[->] (check.east) -- node[above] {是} (branchRight) -- (increment.south);
  \draw[->] (increment.north) -- (sum.south);
  \draw[->] (sum.north) -- (2.55,8.35) -- (loopJoin);
  \draw[->] (check.south) -- node[right] {否} (output.north);
\end{tikzpicture}
